\section{动态比值误差与共模瞬态泄漏}
\label{sec:dynamic}

第~\ref{sec:model-plane} 节的缓冲前端将 $v_P=IR$ 和 $v_N=IR_s$ 分别送入两条独立的全正值元件 RC 梯形网络，端接负载为恒定的非负电阻负载。电阻变化改变信号源幅度，各梯形网络仍为稳定的线性时不变系统。下文的卷积包含完整的输入历史，或采用与其等效且相容的初始状态，并要求参考通道输出非零。记 $H_j=1/D_j$、$G_j=H_j/H_j(0)=D_j(0)/D_j$，以 $g_j$ 表示 $G_j$ 的冲激响应。测得输出电压 $u_P,u_N$ 后，经直流增益校准得到
\begin{equation}\label{eq:ratio-readout}
\hat R=R_s\frac{H_N(0)}{H_P(0)}\frac{u_P}{u_N}
=\frac{(g_P\ast IR)(t)}{(g_N\ast I)(t)}.
\end{equation}

\subsection{匹配通道：精确的加权平均恒等式}
\label{sec:dynamic-matched}

当 $G_P=G_N=G$ 时，减去瞬时电阻可得
\begin{equation}\label{eq:weighted-average}
\hat R(t)-R(t)=\frac{\displaystyle\int_0^\infty g(u)I(t-u)[R(t-u)-R(t)]\,\mathrm du}
{(g\ast I)(t)}.
\end{equation}
对于正值激励，读数是电阻历史值的加权平均。当电阻恒为 $R_0$ 时，对于所有使分母非零的激励波形，读数均等于 $R_0$。当电阻变化时，激励仍参与确定权重。例如，取 $I=I_0e^{\alpha t}$、$R=R_0e^{\beta t}$，其中 $\alpha,\beta\ge0$，且完整历史均满足相应的指数关系，则有
\begin{equation}\label{eq:dynamic-exp}
\hat R(t)=R(t)\frac{G(\alpha+\beta)}{G(\alpha)}.
\end{equation}
即使通道完全相同，改变 $\alpha$ 通常仍会改变读数（图~\ref{fig:dynamic-weighting}）。

对于全正值元件的对地 RC 梯形网络，节点状态矩阵的非对角元素均非负，输入和输出向量也均非负。因此，其矩阵指数给出 $g(t)\ge0$。部分分式展开的留数符号交替，并不意味着冲激响应存在负值区段。于是，
\begin{equation}\label{eq:dynamic-moment}
M_1=\int_0^\infty u|g(u)|\,\mathrm du
=\int_0^\infty ug(u)\,\mathrm du=\tau,
\qquad \tau=T\frac{D_x'(0)}{D(0)},\quad T=R_*C_*.
\end{equation}
若在相关输入历史内满足 $|\dot R|\le L_R$ 和 $|I|\le I_{\max}$，则有
\begin{equation}\label{eq:lag-bound}
|\hat R-R|\le\frac{L_R I_{\max}\tau}{|(g\ast I)(t)|}.
\end{equation}
若在该段完整历史内始终有 $0<I_{\min}\le I\le I_{\max}$，则上界为 $L_R(I_{\max}/I_{\min})\tau$。若从有限时刻启动，则需要使用实际分母。

\subsection{失配通道：电流变化率的贡献}
\label{sec:dynamic-mismatch}

归一化响应核不同时，误差的精确分解为
\begin{equation}\label{eq:dynamic-decomposition}
\hat R-R=\frac{\displaystyle\int_0^\infty g_P(u)I(t-u)[R(t-u)-R(t)]\,\mathrm du}{(g_N\ast I)(t)}
+R\frac{((g_P-g_N)\ast I)(t)}{(g_N\ast I)(t)}.
\end{equation}
两项均采用参考通道的分母。当 $I,R$ 缓慢变化，且分母的绝对值具有大于零的下界时，展开各卷积可得
\begin{equation}\label{eq:mismatch-dynamic}
\hat R-R\simeq-\tau_P\dot R-R\frac{\dot I}{I}\Delta\tau,
\qquad \Delta\tau=\tau_P-\tau_N.
\end{equation}
其中 $\tau_j=T(A_{j,1}+gB_{j,1})/(1+gB_j(0))$。可用 $\mu_{2,j}=\int u^2g_j(u)\,\mathrm du$ 控制展开余项：对于二阶可导输入 $f$，
\begin{equation}\label{eq:dynamic-remainder}
|(g_j\ast f)-f+\tau_j\dot f|
\le\tfrac12\mu_{2,j}\sup_{v\le t}|f''(v)|.
\end{equation}
先分别将该界用于 $f=IR$ 和 $f=I$，再展开二者的比值。对于快速变化的边沿，需要使用精确表达式。参考通道较慢时，$\Delta\tau<0$；在电阻恒定、电流上升时，会产生正误差（图~\ref{fig:dynamic-slew}）。

\subsection{恒定电阻下的匹配判据}
\label{sec:constant-resistance-match}

校准后的读数对所有允许的电流历史均等于任意恒定电阻 $R_0$ 的充要条件为 $g_P=g_N$。充分性由式~\eqref{eq:ratio-readout} 得到。对于必要性，若每个探测电流均满足 $(g_P-g_N)\ast I=0$，则两个响应核必须相等。在正值恒定电流附近施加小扰动，已足以作为探测信号。等价地，
\begin{equation}\label{eq:dynamic-normalized-match}
G_P=G_N\quad\Longleftrightarrow\quad
\frac{D_P}{D_P(0)}=\frac{D_N}{D_N(0)}.
\end{equation}
未归一化分母相等，即 $D_P=D_N$，还要求两个直流增益相等。只有同时满足该要求，表~\ref{tab:load-dimensions} 中按未归一化传递函数计算的匹配族维数才适用。当电阻变化时，匹配的响应核给出式~\eqref{eq:weighted-average}，其中的权重仍取决于激励。

\subsection{负载改变后的共模瞬态泄漏}
\label{sec:cm-pulse}

对于在 $g_0$ 下匹配、在 $g$ 下观测的支路，记 $\Delta B=B_P-B_N$，由式~\eqref{eq:general-load-shift} 得到 $F=G_{DC}=-(g-g_0)\Delta B/(D_PD_N)$。由于 $F(0)=0$ 且极点稳定，幅度为 $V_c$ 的共模阶跃产生衰减的差分响应。以 $p_k$ 表示采用物理频率单位、约分后仍保留的互异极点，以 $a_k$ 表示 $F(s)/s$ 的留数，则有
\begin{equation}\label{eq:cm-pulse}
y(t)=V_c\sum_k a_ke^{p_kt},\qquad
\int_0^\infty y(t)^2\,\mathrm dt
=V_c^2\sum_{k,j}\frac{a_ka_j}{-(p_k+p_j)}.
\end{equation}
若改以 $r_k$ 表示 $F$ 的留数，则 $a_k=r_k/p_k$。冲激响应能量的求和项采用 $r_kr_j$，与上述阶跃响应能量是不同的量。此外，$\sup|y|\le |V_c|\|h_{DC}\|_1$，其中 $h_{DC}=\mathcal L^{-1}F$。存在重极点时，需要使用相应的多项式与指数函数乘积项。

表~\ref{tab:branches} 中的 P/T 支路在 $g_0=0.1$ 下匹配，在 $g=0.2$ 下观测时，P 减去 T 的传递函数为
\begin{equation}\label{eq:dynamic-pt}
F(x)=\frac{0.1x}{(x^2+3.2x+1.4)(x^2+3.3x+1.4)}.
\end{equation}
其无量纲极点为 $-2.8,-2.677033,-0.522967,-0.5$。取 $T=10~\mu\mathrm s$，单位阶跃响应在 $28.20333~\mu\mathrm s$ 达到峰值 $0.009197601$（图~\ref{fig:dynamic-pulse}）；约分后保留的最慢极点给出 $20~\mu\mathrm s$ 的渐近时间常数。阶跃响应能量为 $3.17497\times10^{-9}V_c^2~\mathrm s$；要确定有限的稳定时间，还需指定容差。




\subsection{瞬态验证}
\label{sec:dynamic-verif}

指数激励的评估窗口为 $2$--$3$~ms，最大步长为 $0.1~\mu\mathrm s$；正弦激励的窗口为 $20$--$40$~ms，最大步长为 $1~\mu\mathrm s$。细化后的最大步长分别为 $0.05~\mu\mathrm s$ 和 $0.5~\mu\mathrm s$。

六个基准原生瞬态算例及两个步长细化算例检验了三级和四级归一化元件值均为 1 的梯形网络，采用 $R_*=1~\mathrm{k\Omega}$、$C_*=10~\mathrm{nF}$ 及 $10~\mathrm{k\Omega}$ 端接负载。四个匹配通道算例采用 $R=10~\mathrm{k\Omega}\,e^{200t}$，激励分别为恒定电流和与 $e^{1000t}$ 成正比的电流。在 $2.5~\mathrm{ms}$ 时，改变激励使匹配通道读数在 $N=3$ 时变化 $5.11377~\Omega$，在 $N=4$ 时变化 $11.94461~\Omega$，与式~\eqref{eq:dynamic-exp} 一致。

两个恒定电阻算例采用 $R=100~\mathrm{k\Omega}$、$I=8~\mu\mathrm A[1+0.2\sin(2\pi100t)]$，并将参考通道末级电容增大 $10\%$。两者分别得到 $\Delta\tau=-2.30769~\mu\mathrm s$ 和 $-2.85714~\mu\mathrm s$。保留符号的一阶预测与精确响应之间的相对均方根偏差分别为 $5.59\%$ 和 $8.63\%$。所有算例中，初始过渡结束后，仿真比值读数与精确值的最大偏差小于 $1.23\times10^{-6}~\Omega$。两个 $N=4$ 的时间步长加密算例使读数变化小于 $1.4\times10^{-6}~\Omega$。

\begin{figure}[htbp]
\centering
\includegraphics[width=0.67\linewidth]{figures/Fig_08_dynamic_weighting.pdf}
\caption{匹配三阶、四阶梯形支路中电阻比值对电流历程的依赖。纵轴为同一电阻 $R(t)=10\,\mathrm{k\Omega}\exp(200t)$ 在两种电流 $I_\alpha(t)=1\,\mathrm{\mu A}\exp(\alpha t)$ 下的比值差 $\widehat R_{\alpha=1000}(t)-\widehat R_{\alpha=0}(t)$，其中 $t$ 以秒计，$\alpha$ 的单位为 $\mathrm{s}^{-1}$。各串联电阻均为 1~k$\Omega$，各对地电容均为 10~nF，两条负载均为 10~k$\Omega$。实线为指数输入的解析响应，空心标记为原生 LTspice 比值插值后的差值。}
\label{fig:dynamic-weighting}
\end{figure}

\begin{figure}[htbp]
\centering
\includegraphics[width=0.67\linewidth]{figures/Fig_09_dynamic_slew.pdf}
\caption{四阶梯形支路对在恒定电阻 $R=100$~k$\Omega$ 和电流 $I(t)=8\,\mathrm{\mu A}[1+0.2\sin(2\pi\,100t)]$ 下的电阻比值误差，其中 $t$ 以秒计。两条支路均采用 1~k$\Omega$ 电阻、10~nF 电容和 10~k$\Omega$ 负载，参考支路末级对地电容增至 11~nF。黑线为精确稳态相量比值，红线为一阶矩预测 $-R\Delta\tau\dot I/I$，空心圆为原生 LTspice 数据。其中 $\Delta\tau=\tau_P-\tau_N=-2.85714~\mathrm{\mu s}$。}
\label{fig:dynamic-slew}
\end{figure}

\begin{figure}[htbp]
\centering
\includegraphics[width=0.67\linewidth]{figures/Fig_10_dynamic_pulse.pdf}
\caption{P/T 支路对在 5~k$\Omega$ 对地负载下的正向共模转换脉冲 $v_{oP}-v_{oT}$。共模输入从 1.0005~ms 开始，在 10~ns 内由 0 升至 1~V；横轴为相对该上升沿起点的时间。实线为有限斜坡输入的解析响应，空心圆为逐点读取原生 LTspice 文件得到的输出差值。在 10~k$\Omega$ 负载下匹配的设计，负载改为 5~k$\Omega$ 后产生瞬态输出差。}
\label{fig:dynamic-pulse}
\end{figure}

\FloatBarrier
